合并后的600150是一个覆盖多船型、多基地、修船与部分装备的综合平台，但资本市场只有一个可交易母体。七家主要船厂的能力和盈利差异有助于解释合并结果，不能转化为七个独立估值单元。投资问题因此从“有哪些船厂”转向“这些船厂能否在统一股本下改善交付、毛利和现金”。

\section{E09｜七家主要船厂与产品能力}

产品矩阵覆盖常规商船、高端气体/绿色船、邮轮、海工、特种与修船，使平台能够承接不同需求周期；但船型级合同额、排程、利用率和毛利未披露，能力只能作为收入和交付的方向证据。

\begin{exhibitbox}[主要船厂能力地图]
\small
\begin{tabularx}{\exhibitboxwidth}{L{2.3cm}L{5.0cm}L{3.0cm}X}
\toprule
\textbf{经营实体} & \textbf{主要产品/能力} & \textbf{合并状态} & \textbf{关键边界} \\
\midrule
江南造船 & 大型集装箱、LPG/氨/乙烷/乙烯等气体船、PCTC、特种船 & 并表组成 & 高端能力成立；逐船经济性未披露 \\
大连造船 & VLCC/油轮、集装箱、散货、化学品、海工及特种 & 吸收合并后并表 & 多产品执行与整合需验证 \\
外高桥造船 & 大型邮轮、散货、集装箱、油轮、PCTC、FPSO & 并表组成 & 邮轮学习曲线不能独立资本化 \\
广船国际 & 油化船、Aframax、PCTC、客滚、半潜及特种 & 并表组成 & 船型结构较好，持续性待交付验证 \\
武昌造船 & 公务科研、支线集装箱、多用途、特种与海工 & 吸收合并后并表 & 军品经济性未披露 \\
中船澄西 & 中型散货、支线集装箱、MR油轮、特种运输、修船 & 并表组成 & 修船毛利与坞期利用率未披露 \\
北海造船 & 大型散货/VLOC、VLCC、集装箱、养殖与修船 & 吸收合并后并表 & 大型船交付与周期敏感度较高 \\
\bottomrule
\end{tabularx}
\sourcenote{CF-03、CF-12；经营实体均为600150合并组成。}
\end{exhibitbox}

该矩阵的投资含义在于风险分散与执行复杂度并存。多船型、多基地降低单一产品需求波动，却增加排程、供应链、内部交易和治理协同难度。只有合并费用、交付效率和营运资金出现可量化改善，平台宽度才转化为每股价值。

\section{E10｜子公司盈利差异是观察信号，不是SOTP}

2025年主要子公司净利率代理差异显著：外高桥和北海较高，江南与大连较低。这可能反映船型、交付节点、内部交易、少数股东和项目生命周期，而不是永久效率排序。删除内部交易和统一净资产口径前，不能把子公司利润直接相加或套独立倍数。

\begin{exhibitbox}[2025年主要子公司经营观察]
\small
\begin{tabularx}{\exhibitboxwidth}{L{2.4cm}R{2.3cm}R{2.3cm}R{2.1cm}X}
\toprule
\textbf{子公司} & \textbf{收入} & \textbf{净利润} & \textbf{净利率代理} & \textbf{研究解读} \\
\midrule
江南造船 & 421.50 & 9.53 & 2.3\% & 高端能力强，项目阶段与结构影响利润 \\
大连造船 & 321.45 & 13.24 & 4.1\% & 油轮/海工平台重要，整合后效率待验 \\
外高桥造船 & 164.69 & 25.04 & 15.2\% & 当期盈利突出，不永久化外推 \\
广船国际 & 201.88 & 20.01 & 9.9\% & PCTC/油化船结构较优 \\
中船澄西 & 87.22 & 8.65 & 9.9\% & 中型商船和修船平衡周期 \\
北海造船 & 116.52 & 16.65 & 14.3\% & 大型散货/VLOC兑现较强，持续性待验 \\
\bottomrule
\end{tabularx}
\sourcenote{CF-03。金额单位：亿元；净利率为分析代理，不代表公司披露分部利润率。}
\end{exhibitbox}

对子公司盈利差异的“所以呢”是把它变成监测问题：高盈利船厂的项目结构能否扩散，低盈利船厂是否因交付时点暂时承压，整合是否降低费用与返工。它们影响合并毛利敏感性，却不增加估值母体数量。

\section{上市体内、参股与集团体外}

中国动力约20.19\%权益属于参股，不是600150并表收入；沪东中华属于CSSC group姐妹单位，不在年报列示的七家主要船厂中。集团层面的LNG能力、订单和全球排名只能作竞争背景，除非600150交易所公告明确资产或合同归属。\srcid{CF-03, CF-12}

\begin{exhibitbox}[唯一估值母体与观察节点]
\small
\begin{tabularx}{\exhibitboxwidth}{L{3.2cm}L{3.0cm}L{3.2cm}X}
\toprule
\textbf{节点} & \textbf{资本关系} & \textbf{研究用途} & \textbf{估值信用} \\
\midrule
600150.SH & 唯一上市估值母体 & 合并盈利、现金与目标价 & 完整合并情景 \\
七家主要船厂 & consolidated components & 交付、产品与经营质量解释 & 不单独估值，不重复加总 \\
中国动力 & 参股约20.19\% & 船用动力边界观察 & 不并入营业收入，不给供应商直达信用 \\
沪东中华 & 集团姐妹单位/上市体外 & LNG能力与竞争背景 & 订单、利润、潜在注入均为0 \\
海外同业 & 独立上市公司 & 周期、效率和高端能力对标 & 本案不做可比倍数 \\
\bottomrule
\end{tabularx}
\sourcenote{CF-03、CF-12；海外同业公开能力资料仅作背景。}
\end{exhibitbox}

这套边界保护了估值算术：所有经营成果最终进入600150合并归母和75.2562亿股分母；任何子公司、参股或集团外资产都不能再加一次。若未来发生资产注入，应重新考虑交易价格、新增股本、少数股东与盈利桥，而不是直接把资产标签加入目标价。

\begin{keyinsight}[公司能力的可投资含义]
合并平台的护城河是多基地船位、复杂产品实绩、工程组织与高额在手订单；其估值约束是透明度、交付复杂度和现金占用。行动上应监测合并效率，而不是按单个明星船厂或集团高端项目追逐估值溢价。
\end{keyinsight}
